\subsection*{NT-E10. The degree-one endpoint reduces to finitely many CM norm forms}
Let $E/k$ have complex multiplication by the imaginary quadratic
field $K$, let $F=kK$, and put $D=|\operatorname{disc}K|$.
Fix $c>0$. This statement is unconditional. Consider good prime
ideals $\mathfrak p$ of $k$ of absolute residue degree one, so that
$N\mathfrak p=p$ is a rational prime, and fix either sign
$a=\pm a_{\mathfrak p}$. Write
\[
 \delta=2\sqrt p-a,\qquad B=4p-a^2 .
\]
Outside the finite ramification set of $F/k$, a prime satisfying
$0\le\delta<c/\sqrt p$ and $p>c/2$ must split in $F/k$.
At every such split prime there is a nonzero integer $b$ with
\[
 B=Db^2,\qquad
 \pi=\frac{a+b\sqrt{\operatorname{disc}K}}2\in\mathcal O_K,\qquad
 N_{K/\mathbb Q}\pi=p.                                  \tag{NT-E10a}
\]
For each fixed $c$, at all sufficiently large split good
degree-one primes the endpoint event is precisely
\[
 0\le2\sqrt p-a<\frac c{\sqrt p}
 \quad\Longleftrightarrow\quad
 a>0\ \hbox{and}\quad 0<Db^2<4c.                        \tag{NT-E10b}
\]
The equality $Db^2=4c$ never qualifies. Consequently, if $4c\le D$,
there are only finitely many degree-one good primes satisfying
either endpoint event. Outside the ramification set of $F/k$,
every possible exception then has $p\le c/2$.

\paragraph{Complete algebraic reduction.}
At a split good prime, all geometric CM endomorphisms specialize
to endomorphisms defined over $\mathbb F_p$. Their rational algebra
$K$ embeds in $\operatorname{End}^0(E_{\overline{\mathbb F}_p})$.
Here specialization is faithful: for $\ell\ne p$, good reduction
identifies the Tate modules, the action of an endomorphism on them
is compatible with specialization, and a nonzero endomorphism
cannot kill all $\ell$-power torsion. Indeed a nonzero endomorphism
of an elliptic curve is an isogeny with finite kernel.

The centralizer of this $K$ inside the rational endomorphism
algebra is $K$ itself. To check this without a classification of
the reduction, use the faithful action
\[
 \operatorname{End}^0(E_{\overline{\mathbb F}_p})
      \otimes_{\mathbb Q}\mathbb Q_\ell
       \hookrightarrow\operatorname{End}_{\mathbb Q_\ell}(V_\ell E)
\]
on the two-dimensional Tate module. The usual integral
injectivity follows from the fact that an endomorphism killing
all $\ell^n$-torsion factors through $[\ell^n]$; passing to the
$\ell$-adic completion and then inverting $\ell$ gives the displayed
map. The centralizer of
$K\otimes\mathbb Q_\ell$ there has dimension two: if the quadratic
algebra is a field the module has rank one over it, and if it is
split its two faithful eigenspaces each have dimension one.
Thus the rational centralizer has dimension at most two and
already contains the two-dimensional algebra $K$.

The $p$-power Frobenius endomorphism commutes with the specialized
CM endomorphisms, because they are defined over $\mathbb F_p$.
It therefore belongs to $K$. Its characteristic polynomial is
$X^2-a_{\mathfrak p}X+p$, so it is an algebraic integer of norm
$p$ and trace $a_{\mathfrak p}$. For the opposite sign replace
Frobenius by its negative, which still belongs to $\mathcal O_K$
and has norm $p$. Writing an algebraic integer in a quadratic
field in the usual integral basis gives
\[
 \pi=(a+b\sqrt{\operatorname{disc}K})/2,\qquad
 a\equiv b\,\operatorname{disc}K\pmod2,\qquad
 4p=a^2+Db^2.
                                                               \tag{NT-E10c}
\]
The integer $b$ is nonzero: otherwise $\pi$ would be a rational
algebraic integer and $p=\pi^2$ would not be prime.
This proves (NT-E10a). The argument uses only the Frobenius
characteristic polynomial, good-reduction specialization and
faithful Tate-module actions; no CM reduction-density theorem
or GRH is required.

\paragraph{Complete analytic inequality at the boundary.}
For a rational prime $p$, $2\sqrt p$ is not integral, so the Hasse
bound is strict and $B>0$. The identity
\[
 \delta\sqrt p=\frac{B}{2+a/\sqrt p}                    \tag{NT-E10d}
\]
holds, with a positive denominator smaller than $4$. Hence
$\delta\sqrt p<c$ implies $B<4c$ strictly.
If $p>c/2$, the same inequality forces $a>0$.
Conversely, suppose $a>0$, $p>c/2$, and $B<4c$.
Both $a$ and $2\sqrt p-c/\sqrt p$ are then positive.
Squaring the desired strict inequality shows the exact equivalence
\[
 \delta\sqrt p<c
 \quad\Longleftrightarrow\quad
 B<4c-\frac{c^2}{p}
 \quad\Longleftrightarrow\quad
 p>\frac{c^2}{4c-B}.                                   \tag{NT-E10e}
\]
There are only finitely many possibilities $B=Db^2<4c$.
If this set is nonempty, let $B_{\max}$ be its largest element and
take
\[
 p>\max\left\{\frac c2,\frac{c^2}{4c-B_{\max}}\right\}.
                                                               \tag{NT-E10f}
\]
Then the final inequality in (NT-E10e) holds simultaneously for
all those possibilities. If the set is empty, no split prime
qualifies. This proves (NT-E10b), including exclusion of equality.

For a good prime unramified and nonsplit in $F/k$, the two CM
eigenspaces are interchanged, so $a_{\mathfrak p}=0$, as proved in
NT-E08. Then $\delta\sqrt p=2p$, which fails the endpoint
inequality for $p>c/2$. The remaining ramified primes form a
fixed finite set. If $4c\le D$, (NT-E10a) gives $B\ge D\ge4c$
at every split prime, establishing the stated finiteness and
explicit restriction on the possible unramified exceptions.

\paragraph{What the endpoint gate says.}
The remaining possible rational primes belong to the finite
collection of quadratic norm families
\[
 p=\frac{a^2+Db^2}{4},\qquad
 0<|b|<2\sqrt{c/D},\qquad
 a>0,\qquad a\equiv b\,\operatorname{disc}K\pmod2.        \tag{NT-E10g}
\]
An actual prime of $k$ must still have the indicated
\emph{curve-specific Frobenius}: arbitrary prime values of these
forms are not asserted to have the desired trace for this
particular $E/k$. Thus the reduction neither discards
congruence/sign conditions nor proves infinitude of quadratic
polynomial prime values. It locates exactly the boundary left
open by NT-E09.

Geometrically, the endomorphism field forces the Frobenius norm
equation into the integral lattice $\mathcal O_K$ of the CM
torus. At this endpoint, the transverse lattice coordinate $b$
is bounded, rather than growing with the norm. Work/run:
\texttt{GRH\_NIGHT/number\_theory},
\texttt{cm\_endpoint\_gate.tex}, NT-E10; complete proof above,
using the CM eigenspace and good-reduction framework of NT-E06--08.
